\documentclass[11pt]{article}
\usepackage[a4paper,margin=25mm]{geometry}
\usepackage{amsmath,amssymb,amsthm}
\usepackage{lmodern}
\usepackage[hidelinks]{hyperref}

\newtheorem*{exercise}{Exercise 3.32}
\newcommand{\RW}{\mathsf{RW}}

\title{Probability on Trees and Networks\\Exercise 3.32 and its Author Comment}
\author{Russell Lyons \and Yuval Peres}
\date{Version of 19 August 2026\\Corrected and updated paperback edition}

\begin{document}
\begin{center}
{\large Probability on Trees and Networks\\[2pt]
Exercise 3.32 and its Author Comment}\\[4pt]
Russell Lyons \quad Yuval Peres\\[2pt]
Version of 19 August 2026 -- corrected and updated paperback edition
\end{center}
\small

% BEGIN CONTEXT 3.32
\noindent\textbf{Source and scope (editorial).}
The exercise and the complete Comment below reproduce the English author
text from the fixed \emph{Version of 19 August 2026} of
\href{https://rdlyons.pages.iu.edu/prbtree/book_pb.pdf}{\emph{Probability on Trees and Networks}}:
Exercise 3.32 on physical PDF page 111 (printed page 94), and Comment 3.32
on physical PDF page 623 (printed page 606). Mathematical glyphs are
retypeset from the original PDF. The original book and its separate
\href{https://rdlyons.pages.iu.edu/prbtree/errata2.pdf}{paperback errata}
are retained as source materials; no erratum has been silently applied
to this excerpt.

\medskip
\noindent\textbf{Uncounted same-book context locators (editorial).}
\begin{itemize}
\item Physical page 30 (printed page 13): the definition of a unit flow
on a rooted tree, its conservation at vertices, and the corresponding
tree-flow energy convention.
\item Section 3.2, physical page 98 (printed page 81): the definition of
branching number, its cutset formulation (3.4), and the notation
$\RW_\lambda$ for the walk associated to conductances
$e\mapsto\lambda^{-|e|}$. Theorem 3.5 on physical page 99 (printed
page 82) locates the recurrence/transience threshold at $\operatorname{br}T$.
\item Section 3.3, physical page 100 (printed page 83): the definition of
$N$-subperiodicity through adjacency-preserving injections of descendant
trees, with the image root at distance at most $N$ from the original
root. The notation $T^x$ is introduced on that page.
\item Section 2.5, physical page 55 (printed page 38): the network unit-flow
definition. Theorem 2.11, ``Energy and Transience,'' on physical page 56
(printed page 39) is the finite-energy unit-flow criterion for transience.
\end{itemize}
These are background locators, not additional counted targets or new
proofs. The exercise below does not prescribe a numerical value of
$\operatorname{br}T$. The Comment explicitly uses an external background
estimate from Kuipers and Niederreiter (1974), bottom of page 125; that
estimate is not reproduced or reproved here. This reader is therefore
not presented as self-contained.
% END CONTEXT 3.32
\normalsize

% BEGIN PRIMARY STATEMENT 3.32
\begin{exercise}
Construct a subperiodic tree for which critical homesick random walk
(that is, $\RW_{\operatorname{br}T}$) is transient.
\end{exercise}
% END PRIMARY STATEMENT 3.32

% BEGIN PRIMARY PROOF 3.32
\noindent\textbf{Comment 3.32.}
This seems to be new. Let $\beta$ be an irrational number. Write $\{x\}$
for the fractional part of a real number $x$. For real $x,y$ and
$k\in\mathbb N$, set
\[
f_k(x,y):=
\bigl(\mathbf{1}_{[0,1/2]}(\{x+k\beta\}),
      \mathbf{1}_{[0,1/2]}(\{y+k\beta\})\bigr)
\]
and let $F_n(x,y)$ be the sequence
$\langle f_k(x,y);\ 0\le k<n\rangle$.
Let $T$ be the tree of all finite sequences of the form $F_n(x,y)$ for
$x,y\in\mathbb R$ and $n\in\mathbb N$, together with the null sequence,
which is the root of $T$. Join $F_n(x,y)$ to $F_{n+1}(x,y)$ by an edge.
Then $|T_n|=4n^2$.
(A sequence
$\langle\mathbf{1}_{[0,1/2]}(\{x+k\beta\});\ 0\le k<n\rangle$
changes as $x$ increases from $0$ to $1$ exactly when one of the points
$\{x+k\beta\}$ passes $1/2$ or $1$, so each of these $n$ points
contributes to two changes.)
Let $\mathcal L_2$ be Lebesgue measure on $[0,1]^2$, and for vertices
$w\in T$, set
\[
\theta(w):=\mathcal L_2
\bigl\{(x,y);\ F_{|w|}(x,y)=w\bigr\}.
\]
Now choose $\beta$ not well approximable by rationals -- for example,
set $\beta:=\sqrt5$. Then $\theta$ is approximately uniform on $T_n$:
Given $w\in T_n$, we want to show that the set of $(x,y)$ for which
$F_n(x,y)=w$ has measure at most $c/n^2$. This set is a product of two
intervals. In fact, all $2n$ intervals determined by $\{k\beta\}$ and
$\{k\beta+1/2\}$ for $0\le k<n$ are of length at most $3/n$: see the
argument at the bottom of p.~125 of Kuipers and Niederreiter (1974).
Therefore $\theta$ has finite energy for unit conductances. That is,
simple random walk is transient.
% END PRIMARY PROOF 3.32

\end{document}
